\section{自动奖励与理解信号}
\subsection{图像与语言的目标表示}
自主奖励的核心要求是减少人工工程量。Chelsea 讲述了用\textbf{目标图像}或\textbf{语言描述}自动化奖励的管线：预训练模型提取视觉语义，比较当前状态与目标的 embedding 缩放，从而给出拖拽式奖励。

\begin{figure}[H]
\centering
\includegraphics[width=0.9\textwidth]{slides-images/slide-02.jpg}
\caption{Slide 2 展示目标-当前状态之间的 embedding 距离如何作为即时 reward。}
\end{figure}

\begin{importantbox}{Embedding Distance Reward}
\begin{itemize}
\item 使用 VAE/CLIP 等模型提取 `f(I)`；\\
\item 奖励采用 $r_t = -\|\phi(I_t) - \phi(I_g)\|_2$；\\
\item 可通过 scratch/finetune 弱化域差异，但核心是通过 representation 比较状态。
\end{itemize}
\end{importantbox}

\subsection{视觉-语言 reward 组合}
Slide 03 展示了将自然语言描述与视觉 embedding 结合生成 reward 的流程：LLM 生成 task prompt，VLM 依据 prompt 评估 current frame，最后用 ensemble output 过滤。通过少量 human-in-the-loop 的特殊 cases，系统可以辨别不安全的 reward signal。

\begin{figure}[H]
\centering
\includegraphics[width=0.85\textwidth]{slides-images/slide-03.jpg}
\caption{Slide 3 展示 LLM 与 VLM 结合产生 reward 的体系，并强调 filter 阶段的重要性。}
\end{figure}

\begin{knowledgebox}{LLM + VLM Reward Pipeline}
\begin{itemize}
\item LLM 接受 task prompt，生成 structured description；\\
\item VLM（如 CLIP/BLIP）对当前 frame 打分；\\
\item 通过 confidence filter（如 softmax temperature）控制 reward 强度。
\end{itemize}
\end{knowledgebox}

\subsection{Reward Pipeline 深度解析}
在 slide 03 所示流程基础上，讲者细化了 reward pipeline 的三段：\textbf{signal extraction}（从 sensors 获取图像/语言）、\textbf{affinity scoring}（embedding similarity）、\textbf{confidence gating}（置信度过滤与 ensemble）。这三段必须在每次模型更新后同步改动，才能保证 reward not drift。

\begin{importantbox}{Pipeline 抽象框架}
\begin{itemize}
\item Signal extraction：logging sensors + normalization；\\
\item Affinity scoring：embedding 距离或 classifier probability；\\
\item Confidence gating：用 threshold/ensemble/temporal smoothing filter 掉异常 reward。
\end{itemize}
\end{importantbox}

\subsection{基于成功分类器的自我监督}
讲者借助 few-shot 标注让 robot 学会识别成功状态：收集少量 positive/negative 快照，训练二分类器 $\sigma(s)$，然后该 classifier 输出是 reward 函数。该方式非常适合完成“目标数量未知”的任务。

\begin{knowledgebox}{Success Classifier Pipeline}
\begin{itemize}
\item 收集 50-100 个自定义“成功”帧，生成负样本；\\
\item 训练轻量 CNN/RNN 输出概率；\\
\item 将概率直接用作 reward，或结合时间衰减防止过早收敛。
\end{itemize}
\end{knowledgebox}

\subsection{多模态 reward 过滤}
Chelsea 进一步指出：CLIP/LLM 生成 reward 时，必须加入\textbf{可信度过滤}，避免模型对 irrelevant features（如背景）给予高分。常用做法是 ensemble 多个模态、加入 temporal consistency 正则化。

\begin{warningbox}{不要让 pretrained 模型 hallucination 结果传递给 policy}
Clip-based reward 可能强化背景闪光点。建议：
\begin{itemize}
\item 使用 temporal smoothing 保证 reward 连续；\\
\item 对于多模态判断，采用 majority vote/variance penalty；\\
\item 保留 human-in-the-loop audit 通道，应急 override。
\end{itemize}
\end{warningbox}

\subsection{本章小结}
自动奖励通过图像、语言、模型分类器等多种信号取代传统手工 reward，并强调需要对多模态输出进行置信度控制。
\newpage

